\documentclass[11pt,a4paper]{article}

\usepackage[margin=2.2cm]{geometry}
\usepackage[T1]{fontenc}
\usepackage[utf8]{inputenc}
\usepackage{lmodern}
\usepackage{amsmath,amssymb}
\usepackage{enumitem}
\usepackage{parskip}

\begin{document}

\section*{Suggested Answers --- trial\_exam\_3}

\subsection*{Section 1 --- Multiple Choice (Q1--Q20)}
\begin{itemize}
    \item Q1 B
    \item Q2 A
    \item Q3 A
    \item Q4 C
    \item Q5 A
    \item Q6 A
    \item Q7 B
    \item Q8 B
    \item Q9 C
    \item Q10 A
    \item Q11 A
    \item Q12 C
    \item Q13 C
    \item Q14 A
    \item Q15 B
    \item Q16 B
    \item Q17 B
    \item Q18 B
    \item Q19 A
    \item Q20 A
\end{itemize}

\subsection*{Section 2 --- Short Questions (Q1--Q6)}
\begin{enumerate}[label=\arabic*.]
    \item \textbf{Symbol and $T_s$:} a symbol is the time interval over which a (modulated) signal carries information; $T_s$ is the symbol duration.
    \item \textbf{Symbol rate and $T_s$:} $R_s = 1/T_s$ (symbols per second).
    \item \textbf{ISI:} intersymbol interference occurs when multipath causes the tail of one transmitted symbol to overlap the next symbol at the receiver; this happens when the channel delay spread is significant relative to $T_s$.
    \item \textbf{Spectral efficiency:} $\eta = \frac{R_b}{B}$ and it measures how many information bits per second are carried per Hz of bandwidth.
    \item \textbf{Duplexing:} how a system shares resources between uplink and downlink (time division or frequency division).
    \item \textbf{Resource allocation example:} allocate time and power (or time and spectrum, etc.); challenge: manage interference and QoS across users.
\end{enumerate}

\subsection*{Section 3 --- Long / Applied Questions}

\subsubsection*{Question 1 --- Link Budget and Feasibility}
Given: $f=750$ MHz, $P_t=15$ dBm, $G_t=3$ dBi, $G_r=2$ dBi, $d=1.2$ km, $N=-112$ dBm, $\gamma_{\min}=6$ dB.

\begin{enumerate}[label=\alph*.]
    \item \textbf{Free-space path loss:}
    \[
    L_{\mathrm{FSPL}}=32.44+20\log_{10}(750)+20\log_{10}(1.2)\approx 91.52\ \mathrm{dB}.
    \]
    \item \textbf{Received power:}
    \[
    P_r=P_t+G_t+G_r-L_{\mathrm{FSPL}}=15+3+2-91.52\approx -71.52\ \mathrm{dBm}.
    \]
    \item \textbf{Minimum required received power:}
    \[
    P_{\min}=N+\gamma_{\min}=-112+6=-106\ \mathrm{dBm}.
    \]
    \item \textbf{Feasible?} yes, since $P_r\approx -71.52 > -106$ dBm.
    \item (Optional margin) $\approx 34.48$ dB.
    
\end{enumerate}

\subsubsection*{Question 2 --- Indoor Path Loss}
Given: distance loss $65$ dB, two concrete walls $2\times 7$ dB, door $4$ dB, glass $3$ dB, $P_t=10$ dBm, $G_t=G_r=0$ dBi, sensitivity $-75$ dBm.

\begin{enumerate}[label=\alph*.]
    \item Total path loss:
    \[
    L_{\text{total}}=65+2\cdot 7+4+3=86\ \mathrm{dB}.
    \]
    \item Received power:
    \[
    P_r=P_t-L_{\text{total}}=10-86=-76\ \mathrm{dBm}.
    \]
    \item Link works? With sensitivity $-75$ dBm, $-76$ dBm is $1$ dB below sensitivity, so \textbf{no}.
    \item Two improvements (examples): increase transmit power; use higher-gain antennas / better placement; reduce number of walls (or add relay).
\end{enumerate}

\subsubsection*{Question 3 --- Frequency Reuse}
Given reuse factor $N=3$.

\begin{enumerate}[label=\alph*.]
    \item Fraction of spectrum per cell:
    \[
    \frac{1}{N}=\frac{1}{3}\approx 0.333\ (33.3\%).
    \]
    \item For $N=6$:
    \[
    \frac{1}{6}\approx 0.167\ (16.7\%).
    \]
    \item Impact of $N$:
    \begin{itemize}
        \item Larger $N$ $\Rightarrow$ more separation between co-channel cells $\Rightarrow$ lower co-channel interference.
        \item Smaller $N$ $\Rightarrow$ larger $1/N$ (more bandwidth per cell) and typically higher capacity per cell, but with higher co-channel interference.
    \end{itemize}
    
\end{enumerate}

\subsubsection*{Question 4 --- Multi-hop Routing}
Feasible links: $1\to 2,\ 1\to 3,\ 2\to 3,\ 2\to 4,\ 3\to 4,\ 1\to 4$.

Energies: $E_{12}=2.4,\ E_{13}=4.1,\ E_{23}=1.6,\ E_{24}=2.8,\ E_{34}=1.5,\ E_{14}=7.2$.

\begin{enumerate}[label=\alph*.]
    \item All feasible routes from $1$ to $4$:
    \begin{itemize}
        \item Route A: $1\to 4$
        \item Route B: $1\to 2\to 4$
        \item Route C: $1\to 3\to 4$
        \item Route D: $1\to 2\to 3\to 4$
    \end{itemize}
    \item Total energy per route:
    \[
    E_A=E_{14}=7.2
    \]
    \[
    E_B=E_{12}+E_{24}=2.4+2.8=5.2
    \]
    \[
    E_C=E_{13}+E_{34}=4.1+1.5=5.6
    \]
    \[
    E_D=E_{12}+E_{23}+E_{34}=2.4+1.6+1.5=5.5
    \]
    \item Minimum-energy route: Route B ($1\to 2\to 4$) with $E_{\min}=5.2$.
    \item Advantage: improved coverage/feasibility over long distances. Disadvantage: more delay and overhead (more hops).
\end{enumerate}

\subsubsection*{Question 5 --- Modulation and Throughput}
Bandwidth: $B=150$ kHz $=1.50\times 10^5$ Hz.

\begin{enumerate}[label=\alph*.]
    \item Shannon capacity at SNR $=9$ dB:
    \[
    \mathrm{SNR}_{\text{lin}}=10^{9/10}\approx 7.943,\quad
    C=B\log_2(1+\mathrm{SNR}_{\text{lin}})
    \]
    \[
    C\approx 150000\cdot \log_2(8.943)\approx 4.74\times 10^5\ \text{bit/s}\ (\approx 474.1\ \text{kbps}).
    \]
    \item Highest usable scheme at $9$ dB: scheme-4 (since it requires SNR $\ge 8$ dB). Throughput $\approx 2.0$ bit/s/Hz.
    \[
    R\approx 2.0\cdot B=2.0\cdot 150000=3.00\times 10^5\ \text{bit/s}\ (300\ \text{kbps}).
    \]
    \item SNR drop from scheme-4 to scheme-3 (scheme-4 needs $\ge 8$ dB): starting at $11$ dB,
    \[
    \Delta \approx 11-8=3\ \text{dB}.
    \]
\end{enumerate}

\subsubsection*{Question 6 --- Coverage Planning}
\begin{enumerate}[label=\alph*.]
    \item \textbf{Single-site:} place the base station on the higher point/crest to maximize line-of-sight or diffraction advantage toward both towns.
    \item \textbf{Two-site (base + relay):} place the base station to strongly serve one town, and place the relay so it has a clearer path over/around the hill to reach the shadowed town.
    \item \textbf{Compare:} two-site generally improves coverage and reliability in the shadowed region but increases cost/complexity (extra site, backhaul, and coordination).
    \item \textbf{Most important propagation effects:} shadowing/blockage by the hill, diffraction around obstacles, and free-space path loss with distance (plus multipath near buildings/terrain).
\end{enumerate}

\end{document}

